\documentclass[11pt,oneside]{amsbook}
\usepackage[T1]{fontenc}
\usepackage{amsmath,amssymb,amsthm,mathtools}
\usepackage[colorlinks=true,linkcolor=blue,citecolor=blue,urlcolor=blue]{hyperref}
\newtheorem{theorem}{Theorem}[chapter]
\newtheorem{proposition}[theorem]{Proposition}
\newtheorem{lemma}[theorem]{Lemma}
\theoremstyle{definition}
\newtheorem{definition}[theorem]{Definition}
\theoremstyle{remark}
\newtheorem{remark}[theorem]{Remark}
\setcounter{tocdepth}{1}
\title{Kravchuk Encyclopedia\\ A Working Book Skeleton}

\date{}
\begin{document}
\frontmatter
\maketitle
\chapter*{Editorial Preface}
This editable skeleton follows the eighteen main topics and appendix program in Section 6.1 of the supplied article source, \texttt{KravFeinsKocik.tex}. Additional headings from the supplied archived website, \texttt{Kravchuk.zip}, have been placed within the closest corresponding chapters. The archived introduction now appears as the first numbered chapter; the eighteen topical chapters follow it. The contents preserve the original program rather than impose a new textbook sequence.

Most sections are intentionally empty. A small set of starter passages establishes a consistent symmetric normalization and elementary identities. Those passages are introductory mathematical derivations, not a comprehensive literature review. Recent spectral-asymptotic material is explicitly labeled as an unfinished working program. Empty sections contain source comments for future development, but no artificial prose filler.

The title retains ``Kravchuk'' while the polynomial notation uses the conventional spelling ``Krawtchouk.'' Historical names and original terminology should be checked as chapters are developed. In particular, the operator symbols $A$, $K$, $\Lambda$, and $T$ require definitions from their original source contexts before they are used interchangeably.
\chapter*{Notation and Conventions}
Throughout the starter material, $K_n(j;N)$ is defined by the ordinary generating function, $k_n(N-2j,N)=n!K_n(j;N)$, $\Phi_N=(K_n(j;N))$, $B_N=\operatorname{diag}(\binom Nj)$, and $S_N=\Phi_NB_N$. All matrix indices begin at zero. Physicists' and probabilists' Hermite normalizations must be explicitly distinguished. Binomial coefficients outside their usual integer range vanish.
\tableofcontents
\mainmatter
\chapter{Introduction}
\label{ch:introduction}
\begin{flushright}
\texttt{Omnibus ex nihilo ducendis sufficit unum.}\\
\emph{``To make all things from nothing, unity suffices.''}\\
\textbf{Leibniz}
\end{flushright}
\begin{center}
\textbf{The Kravchuk Encyclopedia}\\
\emph{Ph. Feinsilver and J. Kocik}
\end{center}
So it is written on a medal entitled \emph{Imago Creationis} and designed by Leibniz to \emph{``exhibit to posterity in silver''} his discovery of the binary system.

\medskip
\noindent\textbf{Note on spellings.}
Since the original articles appeared in the French Academy journal \emph{C.R.A.S.}, the French spelling Krawtchouk was used and has been generally adopted. We will use with equal propriety the English spelling Kravchuk, which, in fact, we encourage in the current context.

\medskip
This is the beginnings of a `Kravchuk Encyclopedia'. From the outline presented here, as the project develops, the gaps will gradually be filled in. We hope that you will find lots of interesting material in this work.

% Editorial addition; distinct from the archived introduction.
\section{Mikhail Kravchuk (1892--1942): life and mathematical legacy}
\label{sec:kravchuk-biography}



Mikhail Kravchuk, also known as Mykhailo Pylypovych Kravchuk, was a
Ukrainian mathematician whose work ranged across algebra, differential
and integral equations, probability, and mathematical statistics. He was
born on 27 September 1892 in Chovnitsy, in Volyn, in what is now Ukraine,
and died on 9 March 1942 in a Soviet forced-labour camp in Kolyma.
The biographical outline here draws on the MacTutor History of Mathematics
Archive and the institutional biography maintained by the Institute of
Mathematics of the National Academy of Sciences of Ukraine.\footnote{MacTutor
History of Mathematics Archive, University of St Andrews,
``Mykhailo Pilipovich Krawtchouk,''
\url{https://mathshistory.st-andrews.ac.uk/Biographies/Krawtchouk/};
Institute of Mathematics, National Academy of Sciences of Ukraine,
``Kravchuk Mikhail Pilipovich (1892--1942),''
\url{https://www.imath.kiev.ua/famous?n=kravchuk&lang=en}.
These are secondary biographical accounts, not archival documents.}

\subsection*{Education and scientific career}
Kravchuk's family moved to Lutsk in 1901. He completed his schooling there
with a gold medal in 1910 and entered the University of Kyiv, then
St Vladimir University, graduating in 1914. His early research included
commuting matrices. After the upheavals of the First World War and the
subsequent revolutionary and civil-war years, he combined university
teaching with work in schools. He defended his doctoral dissertation in
1924 and received the title of professor in 1925. His work on matrices,
bilinear forms, and linear transformations was an important part of this
early development.

In 1929 he was elected a full member of the Ukrainian Academy of Sciences.
He taught at the Kyiv Polytechnic Institute and served as head of its
mathematics chair. The Institute of Mathematics identifies him as one of
its founders. His activity was not confined to research: he helped develop
Ukrainian-language mathematical education and terminology, wrote teaching
materials, and supported the training of younger mathematicians.

\subsection*{The polynomials and the scope of his work}
The orthogonal polynomials bearing his name were introduced in 1929.
They are associated with the binomial distribution, providing the
historical starting point for the finite orthogonality and transform
identities developed in this Encyclopedia. Their subsequent uses in
coding theory, finite transforms, and algebraic constructions belong to
the later development of the subject; they should not all be attributed
to Kravchuk himself.

Although these polynomials are now his most familiar mathematical legacy,
his research was considerably broader. It included matrix theory, the
method of moments, differential and integral equations, probability,
statistics, and questions in mathematical education. This wider setting
is particularly relevant here: the interplay between matrices,
orthogonality, and operator methods is not merely an accidental grouping
of topics under one name.

\subsection*{Persecution, death, and rehabilitation}
Kravchuk was arrested in 1938 during Stalin's purges and sentenced on
political charges to twenty years of imprisonment. He was sent to Kolyma,
where he died in 1942, aged forty-nine. MacTutor describes coercion and
threats against his family during the proceedings. He was posthumously
rehabilitated in 1956. His membership in the Ukrainian Academy of Sciences
was restored in 1992. These events are essential to an account of his
career: its abrupt end was a consequence of political repression, not a
natural conclusion to his scientific work.\footnote{For the account of
his persecution, rehabilitation, and restoration of Academy membership,
see the MacTutor biography cited above.}

\subsection*{Names and conventions in this Encyclopedia}
The spellings \emph{Kravchuk}, \emph{Krawtchouk}, and \emph{Krawtschuk}
reflect different transliterations. MacTutor notes that \emph{Krawtchouk}
was the spelling used in his French-language papers. We retain
\emph{Kravchuk} in the Encyclopedia's title and historical discussion,
while respecting \emph{Krawtchouk} in established mathematical terminology,
source titles, and quotations. The first names \emph{Mikhail} and
\emph{Mykhailo} likewise occur in the biographical literature. These
variants refer to the same mathematician.

\medskip
\noindent\emph{Editorial provenance.} This sourced biographical section
is a later addition to the Encyclopedia. It is not part of the archived
2001 introduction, whose wording and attribution are retained separately.


\section*{Starting points}
The original introduction offered the following entry points, now represented within this book:
\begin{itemize}
\item \emph{First View}: see Section~\ref{sec:1-5}.
\item \emph{Krawtchouk Matrices}: see Section~\ref{sec:1-4}.
\item \emph{Kravchuk Polynomials}: see Section~\ref{sec:1-2}.
\item \emph{Twistons}: see Section~\ref{sec:9-6}.
\item \emph{Other Examples}: reserved for future additions to the introductory material.
\end{itemize}
For transforms, see Chapter~\ref{ch:14}. The complete table of contents replaces the original website's \emph{To Contents} link.

\section*{Archival note}
This chapter transcribes the substantive text of \texttt{Kravchuk/Introduction.html} in the supplied website archive. The displayed author attribution and the historical wording above are retained from that source; they do not assign authorship of later editorial additions. Its modification footer records July 28, 2001. The historical claims and quotation are reproduced as source text and have not been independently verified for this skeleton.

The original page also linked to an interactive K-transform form and a bibliographic search database. Those were website services rather than book text; they are recorded here as archival resources, not represented as functioning services. Navigation has been adapted to book cross-references.

\chapter{Pascal's Triangle}
\label{ch:1}
\section{Pascal's triangle and binomial identities}
\label{sec:1-1}
Choosing an $n$-element subset of $\{1,\ldots,N\}$ according to whether it contains the last element proves
\begin{equation}\label{eq:pascal-subsets}
\binom Nn=\binom{N-1}n+\binom{N-1}{n-1}.
\end{equation}
We use $\binom Nn=0$ outside $0\le n\le N$. Expanding a product of $N$ factors gives
\[
(1+v)^N=\sum_{n=0}^N\binom Nn v^n.
\]
Taking coefficients in $(1+v)^{a+b}=(1+v)^a(1+v)^b$ gives Vandermonde's identity
\[
\binom{a+b}n=\sum_r\binom ar\binom b{n-r}.
\]
Replacing the second factor by $(1-v)^b$ yields the signed convolution defining $K_n(b;a+b)$ in Section~\ref{sec:1-2}. Thus the first column of the Krawtchouk array is Pascal's row, and its other columns are signed subset counts, not probability distributions.
For example, at $N=3$, the columns for $j=0$ and $j=1$ are respectively $(1,3,3,1)^T$ and $(1,1,-1,-1)^T$.

\section{Krawtchouk polynomials}
\label{sec:1-2}
For $0\leq n,j\leq N$, define the symmetric Krawtchouk polynomial values by
\begin{equation}\label{eq:gen}
(1+v)^{N-j}(1-v)^j=\sum_{n=0}^N K_n(j;N)v^n.
\end{equation}
Hence
\[
K_n(j;N)=\sum_r(-1)^r\binom jr\binom{N-j}{n-r}.
\]
Binomial coefficients outside their usual range are understood to vanish.
In centered coordinates $x=N-2j$, write
\[
(1+v)^{(N+x)/2}(1-v)^{(N-x)/2}
=\sum_{n=0}^N\frac{v^n}{n!}k_n(x,N).
\]
Thus $k_n(N-2j,N)=n!K_n(j;N)$.

\section{Generating functions}
\label{sec:1-3}
Write $x=N-2j$ and use the generating function in \eqref{eq:gen}. Logarithmic differentiation gives the polynomial identity
\[
(1-v^2)\frac{\partial G_{N,j}}{\partial v}=(x-Nv)G_{N,j}.
\]
Coefficient comparison, with $K_{-1}=0$, yields
\[
(n+1)K_{n+1}(j;N)=xK_n(j;N)-(N-n+1)K_{n-1}(j;N).
\]
In particular,
\[
K_0=1,\qquad K_1=x,\qquad K_2=\frac{x^2-N}{2},\qquad
K_3=\frac{x^3-(3N-2)x}{6}.
\]
The factorial-scaled polynomials are $k_n=n!K_n$, so their leading coefficient as polynomials in $x$ is one. This is why factorials must be retained when passing between ordinary and exponential generating functions.
At the endpoints,
\[
K_n(0;N)=\binom Nn,\qquad K_n(N;N)=(-1)^n\binom Nn.
\]
Moreover, $G_{N,N-j}(v)=G_{N,j}(-v)$, giving $K_n(N-j;N)=(-1)^nK_n(j;N)$. These are fixed-order identities, distinct from the cross-order construction of Section~\ref{sec:1-5}.

\section{Krawtchouk matrices}
\label{sec:1-4}
Set $\Phi_N=(K_n(j;N))_{n,j=0}^N$ and
$B_N=\operatorname{diag}(\binom N0,\ldots,\binom NN)$. Binomial duality is
\[
\binom Nj K_n(j;N)=\binom Nn K_j(n;N).
\]
Consequently $S_N=\Phi_NB_N$ is symmetric. The symmetric transform satisfies
\[
\Phi_N^2=2^NI.
\]
For orientation,
\[
\Phi_1=\begin{pmatrix}1&1\\1&-1\end{pmatrix},\qquad
\Phi_2=\begin{pmatrix}1&1&1\\2&0&-2\\1&-1&1\end{pmatrix}.
\]

\section{First view and cross properties}
\label{sec:1-5}
Write $G_{N,j}(v)=(1+v)^{N-j}(1-v)^j$, as in \eqref{eq:gen}.
Multiplying by $1+v$ or $1-v$ and comparing coefficients gives
\begin{align}
K_n(j;N+1)&=K_n(j;N)+K_{n-1}(j;N),\label{eq:cross-plus}\\
K_n(j+1;N+1)&=K_n(j;N)-K_{n-1}(j;N).\label{eq:cross-minus}
\end{align}
Here $0\leq j\leq N$, and coefficients outside $0\leq n\leq N$
vanish at order $N$. Adding and subtracting yields the inverse cross rules
\begin{align*}
K_n(j;N+1)+K_n(j+1;N+1)&=2K_n(j;N),\\
K_n(j;N+1)-K_n(j+1;N+1)&=2K_{n-1}(j;N).
\end{align*}
These are cross-order identities: the order parameter must not be
suppressed when they are used to construct successive tables.
For example, $K_2(1;3)=-1$ and $K_1(1;3)=1$ give
$K_2(1;4)=0$ and $K_2(2;4)=-2$.

There is also a fixed-order local relation. If
\[
\begin{pmatrix}a&b\\c&d\end{pmatrix}
=\begin{pmatrix}K_n(j;N)&K_n(j+1;N)\\
K_{n+1}(j;N)&K_{n+1}(j+1;N)\end{pmatrix},
\]
then $-a+c=b+d$. Indeed,
$(1-v)G_{N,j}=(1+v)G_{N,j+1}$.
Column reversal and row reversal are instead the fixed-order
symmetries proved in Section~\ref{sec:8-1}; these should be distinguished
from the cross-order rules above. Compare \cite[Appendix~C]{KocikFeynman}.


\chapter{Random Walks}
\label{ch:2}
\section{Paths and path integrals}
\label{sec:2-1}
Consider length-$N$ binary paths $w=(w_1,\ldots,w_N)$, where
$w_a\in\{0,1\}$ and $|w|=\sum_aw_a$ is the endpoint coordinate.
Fix $j$ distinguished steps and assign the signed weight
\[
\omega_j(w)=(-1)^{w_1+\cdots+w_j}.
\]
The particular first-$j$ choice is inessential; any fixed subset of $j$
steps gives the same endpoint sums. A path with endpoint $n$ and $r$
right steps among the distinguished steps has weight $(-1)^r$.
There are $\binom jr\binom{N-j}{n-r}$ such paths, so
\begin{equation}\label{eq:signed-path-Krawtchouk}
\sum_{|w|=n}\omega_j(w)
=\sum_r(-1)^r\binom jr\binom{N-j}{n-r}
=K_n(j;N).
\end{equation}
This is a path interpretation of the coefficient formula in
Section~\ref{sec:1-2}, not a new definition. For $j=0$ it counts ordinary
paths and gives $\binom Nn$; for $j>0$, cancellation produces signed
entries of column $j$ of $\Phi_N$. For instance, $N=4$, $j=2$ gives
endpoint sums $(1,0,-2,0,1)$.

The signed Galton-board construction of \cite[Section~6]{KocikFeynman}
is an elementary analogue of a sum over paths. Its weights are not
probabilities. In quantum mechanics one sums amplitudes before taking
an appropriately normalized absolute square; a physical unitary model
requires more structure than this combinatorial identity. The unsigned
walk and the stochastic Ehrenfest transition operator should likewise
not be identified with a unitary evolution.


\subsection{Weighted paths and the two-parameter deformation}
\label{sec:weighted-path-deformation}
Over a commutative ring, define a separate coefficient array
\[
\mathcal K_n(j;N;\alpha,\beta)
=[v^n](1+\alpha v)^{N-j}(1+\beta v)^j.
\]
A consistent path rule gives every left step weight $1$, and a right
step weight $\beta$ at each of $j$ distinguished times and $\alpha$ at
all other times. Thus
\[
\mathcal K_n(j;N;\alpha,\beta)
=\sum_r\binom jr\binom{N-j}{n-r}\beta^r\alpha^{n-r}.
\]
This recovers $K_n(j;N)$ at $(\alpha,\beta)=(1,-1)$ without changing
its normalization. Coefficient comparison gives
\begin{align*}
\mathcal K_n(j;N+1)&=\mathcal K_n(j;N)+\alpha\mathcal K_{n-1}(j;N),\\
\mathcal K_n(j+1;N+1)&=\mathcal K_n(j;N)+\beta\mathcal K_{n-1}(j;N).
\end{align*}
For adjacent entries $\begin{pmatrix}a&b\\c&d\end{pmatrix}$ in this
array, $\beta a+c=\alpha b+d$.
These identities require no division by $\alpha-\beta$ and remain valid
when parameters coincide or are not invertible.

\begin{remark}[A discrepancy in the source's path rule]
The displayed generating function of \cite[Section~6.2]{KocikFeynman}
is the one used here. However, the preceding prose assigns left/right
weights $\alpha,\beta$ only at the $j$ distinguished times, and weights
$1,1$ elsewhere. That different rule produces
$(1+v)^{N-j}(\alpha+\beta v)^j$, not the displayed generating function.
We retain the displayed coefficient array and specify the compatible
rule explicitly. No stochastic normalization, orthogonality, or
involution property is asserted for arbitrary parameters.
\end{remark}

\section{$A$, $K$, and $\Lambda$}
\label{sec:2-2}
For the symmetric walk on $\{0,\ldots,N\}$, use column vectors for
state distributions and the zero-based standard basis $e_j$. Put
\[
A_Ne_j=(N-j)e_{j+1}+je_{j-1},\qquad
\Lambda_N=\operatorname{diag}(N,N-2,\ldots,-N),
\]
with endpoint terms omitted. In this section the matrix symbol $K$
is explicitly defined to mean $\Phi_N=(K_n(j;N))_{n,j=0}^N$.
It is not a single polynomial $K_n$ or the symmetric matrix
$S_N=\Phi_NB_N$.

\subsection{Intertwining and diagonalization}
Writing $G_j(v)=(1+v)^{N-j}(1-v)^j$, differentiation gives
\[
(1-v^2)G_j'(v)=(N-2j-Nv)G_j(v).
\]
Equating coefficients yields
\[
(N-n+1)K_{n-1}(j;N)+(n+1)K_{n+1}(j;N)
=(N-2j)K_n(j;N).
\]
The spatial recurrence is
\[
jK_n(j-1;N)+(N-j)K_n(j+1;N)=(N-2n)K_n(j;N).
\]
Absent endpoint terms are zero. In matrix form these identities read
\[
A_N\Phi_N=\Phi_N\Lambda_N,\qquad
\Phi_NA_N=\Lambda_N\Phi_N.
\]
This is the coefficient proof of the master equation; its symmetric-power
interpretation is given in Section~\ref{sec:8-1}.
Since $\Phi_N^2=2^NI$, we obtain
\[
\Phi_N^{-1}=2^{-N}\Phi_N,\qquad
A_N=2^{-N}\Phi_N\Lambda_N\Phi_N.
\]
Column $j$ of $\Phi_N$ is a right eigenvector with eigenvalue $N-2j$;
row $n$ is a left eigenvector with eigenvalue $N-2n$.

\subsection{Transition probabilities}
For $N\geq1$, $T_N=A_N/N$ is column-stochastic. Hence
$p_{t+1}=T_Np_t$ and, for integers $t\geq0$,
\[
T_N^t=2^{-N}\Phi_N(\Lambda_N/N)^t\Phi_N.
\]
The backward operator on column vectors of function values is $T_N^T$,
not $T_N$; its eigenfunctions are $j\mapsto K_n(j;N)$.
The stationary column distribution is
$2^{-N}(\binom N0,\ldots,\binom NN)^T$, the normalized first column
of $\Phi_N$. The eigenvalue $-1$ of $T_N$ expresses period two;
stationarity does not imply convergence from arbitrary initial data.

\subsection{The cases $N=3$ and $N=5$}
For $N=3$,
\begingroup
\small
\setlength{\arraycolsep}{3pt}
\[
A_{3}=\begin{pmatrix}
0 & 1 & 0 & 0 \\
3 & 0 & 2 & 0 \\
0 & 2 & 0 & 3 \\
0 & 0 & 1 & 0
\end{pmatrix}.
\]
\endgroup
\begingroup
\small
\setlength{\arraycolsep}{3pt}
\[
\Lambda_{3}=\begin{pmatrix}
3 & 0 & 0 & 0 \\
0 & 1 & 0 & 0 \\
0 & 0 & -1 & 0 \\
0 & 0 & 0 & -3
\end{pmatrix}.
\]
\endgroup
The transition eigenvalues are $1,\frac{1}{3},- \frac{1}{3},-1$.
For $N=5$,
\begingroup
\small
\setlength{\arraycolsep}{3pt}
\[
A_{5}=\begin{pmatrix}
0 & 1 & 0 & 0 & 0 & 0 \\
5 & 0 & 2 & 0 & 0 & 0 \\
0 & 4 & 0 & 3 & 0 & 0 \\
0 & 0 & 3 & 0 & 4 & 0 \\
0 & 0 & 0 & 2 & 0 & 5 \\
0 & 0 & 0 & 0 & 1 & 0
\end{pmatrix}.
\]
\endgroup
\begingroup
\small
\setlength{\arraycolsep}{3pt}
\[
\Lambda_{5}=\begin{pmatrix}
5 & 0 & 0 & 0 & 0 & 0 \\
0 & 3 & 0 & 0 & 0 & 0 \\
0 & 0 & 1 & 0 & 0 & 0 \\
0 & 0 & 0 & -1 & 0 & 0 \\
0 & 0 & 0 & 0 & -3 & 0 \\
0 & 0 & 0 & 0 & 0 & -5
\end{pmatrix}.
\]
\endgroup
The transition eigenvalues are $1,\frac{3}{5},\frac{1}{5},- \frac{1}{5},- \frac{3}{5},-1$.
The corresponding $R_N$, $L_N$, $\Phi_N$ and $S_N$ tables are in
Appendix~\ref{app:matrix-tables}. Centered and lattice polynomial tables
are in Appendix~\ref{app:polynomial-tables}.

\section{Nonsymmetric walks}
\label{sec:2-3}
% TODO: Develop this section; add verified references, examples, and exercises.

\noindent\textbf{NONSYMMETRIC WALKS NEEDS TO BE DONE.}


\section{Symmetric Krawtchouk matrices and binomial expectations}
\label{sec:2-4}
% TODO: Develop this section; add verified references, examples, and exercises.

\noindent\textbf{SYMMETRIC KRAWTCHOUK MATRICES AND BINOMIAL EXPECTATIONS NEEDS TO BE DONE.}



\section*{Connection to the supplied article sources}
The supplied source \texttt{KravFeinsKocik.tex}, Section~5, develops the
classical Bernoulli random walk and a multivariate extension.
The supplied \texttt{jkhadamard\_Oct\_08\_00.tex}, Section~5, connects
classical random walks with Krawtchouk matrices. Read these alongside
this chapter using the book's ordinary generating-function normalization;
source conventions should not be silently substituted for it.

\chapter{Urn Model}
\label{ch:3}
\section{Markov chains}
\label{sec:3-1}
In the discrete-time Ehrenfest chain, the state $j$ is the number of balls in one of two urns. Moving one uniformly chosen ball gives
\[
(Pf)(j)=\frac{N-j}{N}f(j+1)+\frac jN f(j-1).
\]
The Krawtchouk difference equation gives
\[
P K_n(\cdot;N)=\left(1-\frac{2n}{N}\right)K_n(\cdot;N).
\]
Here $N\geq1$, and terms outside the state space have zero transition coefficient.

\section{Initial and invariant distributions}
\label{sec:3-2}
The invariant distribution is $w_N(j)=2^{-N}\binom Nj$. Detailed balance follows from
\[
w_N(j)\frac{N-j}{N}=w_N(j+1)\frac{j+1}{N}.
\]
The non-lazy discrete-time chain has period two; invariance should not be confused with convergence to equilibrium from every initial distribution.


\section*{Connection to the supplied article sources}
Section~4 (``Ehrenfest urn model'') of the supplied
\texttt{KravFeinsKocik.tex} belongs to this chapter.
The urn state is the Hamming weight of a binary configuration: flipping
a uniformly selected coordinate moves from $j$ to $j+1$ with probability
$(N-j)/N$ and to $j-1$ with probability $j/N$, for $N\geq1$.
In the book's column-vector convention the transition matrix is
$A_N/N$. The tensor-space realization is discussed in Chapter~\ref{ch:9}.

\chapter{Symmetric Functions and Energy}
\label{ch:4}
\section{Elementary symmetric functions and determinants}
\label{sec:4-1}
For scalars $z_1,\ldots,z_N$, the elementary symmetric functions are
\[
e_n(z_1,\ldots,z_N)=\sum_{1\leq a_1<\cdots<a_n\leq N}
z_{a_1}\cdots z_{a_n},\qquad e_0=1.
\]
Each subset occurs once: repeated indices and ordered tuples are excluded.
For $D=\operatorname{diag}(z_1,\ldots,z_N)$,
\begin{equation}\label{eq:elementary-energy-generating}
\det(I+vD)=\prod_{a=1}^N(1+z_av)
=\sum_{n=0}^N e_n(z_1,\ldots,z_N)v^n.
\end{equation}
If $z_a=\epsilon_a\in\{+1,-1\}$ and precisely $j$ signs are negative,
comparison with the defining generating function \eqref{eq:gen} gives
\begin{equation}\label{eq:sign-energy-Krawtchouk}
e_n(\epsilon_1,\ldots,\epsilon_N)=K_n(j;N).
\end{equation}
Thus the order-$n$ interaction energy depends only on the number of
negative signs, not their positions. The first two nonconstant energies are
\[
e_1=N-2j,\qquad e_2=\frac{(N-2j)^2-N}{2}.
\]
These are ordinary coefficient values, not the factorial-scaled centered
polynomials: $k_n(N-2j,N)=n!e_n$.
The twiston interpretation is given in Section~\ref{sec:9-6}; the signed-path
interpretation in Section~\ref{sec:2-1} gives the same coefficients by a
different combinatorial construction. See also \cite[Sections~1 and~5]{KocikFeynman}.

\section{Traces on Grassmann algebras}
\label{sec:4-2}
Let $V$ have dimension $N$ and let $T:V\to V$ be linear. On the $n$th exterior power define
\[
(\bigwedge^n T)(v_1\wedge\cdots\wedge v_n)
=Tv_1\wedge\cdots\wedge Tv_n.
\]
The trace of this induced map is the sum of the principal $n\times n$ minors of $T$. Consequently,
\begin{equation}\label{eq:exterior-trace-determinant}
\det(I+vT)=\sum_{n=0}^N v^n\operatorname{tr}(\bigwedge^n T).
\end{equation}
Indeed, in an exterior basis indexed by subsets, the diagonal entry indexed by a subset is precisely its principal minor; the determinant expansion gives the same coefficients. This argument does not require diagonalizability.
For a sign operator $D$ with $N-j$ eigenvalues $+1$ and $j$ eigenvalues $-1$, the degree-$n$ trace is therefore
\[
\operatorname{tr}(\bigwedge^n D)=K_n(j;N).
\]
For the full Grassmann algebra $\bigwedge V=\bigoplus_n\bigwedge^n V$, ordinary trace and supertrace are respectively
\[
\operatorname{tr}(\bigwedge T)=\det(I+T),\qquad
\operatorname{str}(\bigwedge T)=\sum_n(-1)^n\operatorname{tr}(\bigwedge^n T)=\det(I-T).
\]
Here $\bigwedge^0T$ is the identity on the scalar space. These formulas concern multiplicative exterior powers, not the additive Lie-algebra action obtained by applying a generator to one wedge factor at a time. They develop the Grassmann-algebra heading of the supplied source outline from the determinant identity in Section~\ref{sec:4-1}.


\chapter{Martingales}
\label{ch:5}
\section{Iterated integrals}
\label{sec:5-1}
Let $\epsilon_1,\epsilon_2,\ldots$ be independent symmetric signs, $X_m=\sum_{a=1}^m\epsilon_a$, and $\mathcal F_m=\sigma(\epsilon_1,\ldots,\epsilon_m)$. Define
\[
M_m(v)=\prod_{a=1}^m(1+v\epsilon_a)=\sum_{n=0}^m E_n(m)v^n.
\]
Independence and zero mean imply $\mathbb E[M_{m+1}(v)\mid\mathcal F_m]=M_m(v)$. Comparing coefficients gives a martingale in $m$ for every fixed $n$, taking $E_n(m)=0$ for $n>m$. Direct expansion gives
\begin{equation}\label{eq:iterated-sign-sums}
E_n(m)=\sum_{1\le a_1<\cdots<a_n\le m}\epsilon_{a_1}\cdots\epsilon_{a_n}
=\frac{k_n(X_m,m)}{n!}.
\end{equation}
With $E_0(m)=1$, the discrete stochastic integral recursion is
\begin{equation}\label{eq:discrete-iterated-recursion}
E_n(m+1)-E_n(m)=\epsilon_{m+1}E_{n-1}(m),\qquad
E_n(m)=\sum_{a=1}^m E_{n-1}(a-1)\epsilon_a.
\end{equation}
The integrand is evaluated before the new sign is observed. Thus these are iterated sums over strictly increasing times, with no diagonal terms. For example, $E_1(m)=X_m$ and $E_2(m)=(X_m^2-m)/2$.
This develops the iterated-integral interpretation in the probabilistic part of the supplied \texttt{KravFeinsKocik.tex}; all normalizations here follow \eqref{eq:gen}.

\section{Orthogonal functionals}
\label{sec:5-2}
For fixed $N$, independence gives the two-variable generating identity
\[
\mathbb E[M_N(v)M_N(w)]=\prod_{a=1}^N\mathbb E[(1+v\epsilon_a)(1+w\epsilon_a)]=(1+vw)^N.
\]
Consequently,
\begin{equation}\label{eq:martingale-orthogonality}
\mathbb E[E_n(N)E_m(N)]=\delta_{nm}\binom Nn.
\end{equation}
Equivalently, if $J=(N-X_N)/2$ has mass $\pi_N(j)=2^{-N}\binom Nj$, then
\[
\sum_{j=0}^N\pi_N(j)K_n(j;N)K_m(j;N)=\delta_{nm}\binom Nn.
\]
This also follows by expanding into subsets: two distinct subsets leave at least one independent sign to an odd power, so their product has expectation zero. Identical subsets contribute one.
There are $N+1$ nonzero orthogonal functions on the $N+1$-point lattice, so they form a basis of all functions of $J$. Explicitly,
\[
f(j)=\sum_{n=0}^N\widehat f_nK_n(j;N),\qquad
\widehat f_n=\frac{\mathbb E[f(J)K_n(J;N)]}{\binom Nn},\qquad
\mathbb E[f(J)^2]=\sum_{n=0}^N\binom Nn\widehat f_n^2
\]
for real $f$. In factorial normalization the squared norm is $(n!)^2\binom Nn$, not merely $\binom Nn$. The martingale time parameter $N$ should be distinguished from the fixed-order lattice argument $j$.

\section{Krawtchouk polynomials and multinomial distribution}
\label{sec:5-3}
The supplied source extends the product martingale to independent identically distributed finite-valued variables. Let $\mathbb P(Y_a=\xi_r)=p_r$ for $0\le r\le d-1$, with $p_r>0$, $\sum_rp_r=1$, mean $\mu$, and variance $\sigma^2>0$. If $\nu_r$ counts occurrences of $\xi_r$ among $N$ observations, set
\[
\mathcal G_{\boldsymbol\nu}(v)=\prod_{r=0}^{d-1}(1+v(\xi_r-\mu))^{\nu_r}
=\sum_{n=0}^N Q_n(\boldsymbol\nu)v^n.
\]
Then
\[
Q_n(\boldsymbol\nu)=\sum_{a_0+\cdots+a_{d-1}=n}\prod_{r=0}^{d-1}\binom{\nu_r}{a_r}(\xi_r-\mu)^{a_r},
\qquad
Q_n(\boldsymbol\nu+\mathbf e_r)=Q_n(\boldsymbol\nu)+(\xi_r-\mu)Q_{n-1}(\boldsymbol\nu).
\]
The count vector has multinomial mass $N!\prod_r p_r^{\nu_r}/\nu_r!$ on $\sum_r\nu_r=N$. Since a centered increment has mean zero and variance $\sigma^2$,
\[
\mathbb E[\mathcal G_{\boldsymbol\nu}(v)\mathcal G_{\boldsymbol\nu}(w)]
=(1+\sigma^2vw)^N,\qquad
\mathbb E[Q_nQ_m]=\delta_{nm}\binom Nn\sigma^{2n}.
\]
The same conditional-expectation argument as in Section~\ref{sec:5-1} proves the coefficient martingale property as observations are added.

For the two values $+1,-1$ with equal probabilities this reduces exactly to $Q_n=K_n(j;N)$. For $d>2$, these $N+1$ orthogonal functions are generally not a complete basis on the $\binom{N+d-1}{d-1}$ count vectors. The source's scalar-parameter construction is therefore distinguished here from a full multi-index multivariate Krawtchouk system; the latter needs additional independent contrasts and is not asserted by this argument.


\chapter{Lie Algebras and Krawtchouk Polynomials}
\label{ch:6}
\section{$\mathfrak{so}(2,1)$ explained}
\label{sec:6-1}
% TODO: Develop this section; add verified references, examples, and exercises.

\section{$\mathfrak{so}(2,1)$ spinors}
\label{sec:6-2}
% TODO: Develop this section; add verified references, examples, and exercises.

\section{Quaternions and Clifford algebras}
\label{sec:6-3}
% TODO: Develop this section; add verified references, examples, and exercises.

\section{$S$ and $\mathfrak{so}(2,1)$ tensors}
\label{sec:6-4}
% TODO: Develop this section; add verified references, examples, and exercises.

\section{Three-dimensional simple Lie algebras}
\label{sec:6-5}
% TODO: Develop this section; add verified references, examples, and exercises.


\chapter{Lie Groups and Reflections}
\label{ch:7}
\section{Reflections}
\label{sec:7-1}
Let $u$ be a non-null pure split quaternion. The reflection in the
plane $Q$-orthogonal to $u$ is
\[
r_u(v)=-uvu^{-1}.
\]
It sends $u$ to $-u$ and fixes any pure $v$ orthogonal to $u$, since
then $uv=-vu$. Null vectors cannot be used in this formula because
they are not invertible. Composing two such reflections yields a
conjugation action and hence a determinant-one orthogonal map.

For $H=F+G$ one has $H^2=2I$, $H^{-1}=H/2$, and
\[
HFH^{-1}=G,\quad HGH^{-1}=F,\quad HiH^{-1}=-i.
\]
This exchanges the two spacelike axes and reverses the timelike axis.
It preserves $Q$ and has determinant one on the pure space, but lies
outside the time-orientation-preserving component.
See Section~\ref{sec:7-2} for the determinant issue in the spinor space.
See \cite[Section~7]{KocikFeynman}; notation follows the book conventions.

\section{Krawtchouk matrices as group elements}
\label{sec:7-2}
Set $H=F+G=\begin{pmatrix}1&1\\1&-1\end{pmatrix}$. Then
\[
FH=HG,\qquad GH=HF,\qquad H^2=2I.
\]
These elementary identities generate the Krawtchouk intertwining and
reflection symmetries under symmetric powers.

A qualification to the supplied source is necessary:
$\det H=-2$, and $\det(cH)=-2c^2<0$ for every nonzero real scalar $c$.
Thus no real scalar normalization puts $H$ in $SL(2,\mathbb R)$.
Likewise $\det F=\det G=-1$.
The conjugations are valid in $GL(2,\mathbb R)$ (or its projective
quotient); $F,G$ are also traceless Lie algebra elements. These are
distinct statements, and must not be conflated.
Over $\mathbb C$, $(\mathrm{i}/\sqrt2)H$ has determinant one, where
$\mathrm{i}$ denotes the scalar imaginary unit.

The lift to degree-$N$ polynomials is developed once in
Section~\ref{sec:8-1}; orthogonal normalization is treated in
Section~\ref{sec:14-1}. Neither construction requires the real matrix
$H$ to belong to $SL(2,\mathbb R)$.
See \cite[Section~7]{KocikFeynman}; notation follows the book conventions.


\chapter{Representations}
\label{ch:8}
\section{Splitting formula}
\label{sec:8-1}
For $g=\begin{pmatrix}a&b\\c&d\end{pmatrix}\in GL(2,\mathbb R)$ define
on homogeneous degree-$N$ polynomials
\[
(\rho_N(g)p)(x,y)=p(ax+cy,bx+dy).
\]
In the basis $e_j=x^{N-j}y^j$, expansion gives its columns, and
$\rho_N(gh)=\rho_N(g)\rho_N(h)$.
For a Lie algebra element $X$, instead use the derivative
\[
d\rho_N(X)=\left.\frac{d}{dt}\rho_N(I+tX)\right|_{t=0}.
\]
On the full tensor space these two constructions are respectively
\[
g\mapsto g^{\otimes N},\qquad
X\mapsto\sum_{r=1}^N I^{\otimes(r-1)}\otimes X\otimes I^{\otimes(N-r)}.
\]
The group construction is multiplicative; the differential preserves
commutators. Tensoring a Lie algebra element multiplicatively would
not give the second construction.

For the split units and $H=F+G$, direct differentiation gives
\[
d\rho_N(F)e_j=(N-j)e_{j+1}+je_{j-1},\qquad
d\rho_N(G)e_j=(N-2j)e_j.
\]
Meanwhile
\[
\rho_N(H)e_j=(x+y)^{N-j}(x-y)^j
=\sum_{n=0}^N K_n(j;N)e_n.
\]
Thus $d\rho_N(F)=A_N$, $d\rho_N(G)=\Lambda_N$, and
$\rho_N(H)=\Phi_N$. Differentiating the conjugation identity produces
$A_N\Phi_N=\Phi_N\Lambda_N$.

The group lift gives a different consequence. Put
$J_Ne_j=e_{N-j}$ and $D_Ne_j=(-1)^je_j$. Then
\[
\rho_N(F)=J_N,\quad\rho_N(G)=D_N,\quad
J_N\Phi_N=\Phi_ND_N,\quad D_N\Phi_N=\Phi_NJ_N.
\]
Equivalently,
\[
K_{N-n}(j;N)=(-1)^jK_n(j;N),\qquad
K_n(N-j;N)=(-1)^nK_n(j;N).
\]
The master equation and these reflection symmetries therefore arise
from the same two-dimensional identity through different lifts.
See \cite[Section~7]{KocikFeynman}; notation follows the book conventions.

\section{Hilbert space structure}
\label{sec:8-2}
With $w_N(j)=2^{-N}\binom Nj$,
\begin{equation}\label{eq:orth}
\sum_{j=0}^Nw_N(j)K_n(j;N)K_m(j;N)=\binom Nn\delta_{nm}.
\end{equation}
Indeed, the weighted product of two generating functions equals $(1+uv)^N$, and coefficient comparison proves \eqref{eq:orth}. Thus
\[
\psi_{n,N}(j)=\sqrt{\frac{w_N(j)}{\binom Nn}}K_n(j;N)
\]
is an orthonormal basis for the counting-measure space on $\{0,\ldots,N\}$.


\subsection{Column vectors, covectors, and the binomial metric}
\label{sec:binomial-dual-metrics}
For state columns $u,v$ define
\[
\langle u,v\rangle_{B^{-1}}=u^TB_N^{-1}v.
\]
For row covectors $\alpha,\beta$, the dual inner product is
$\alpha B_N\beta^T$. The two weights are inverse to one another,
not competing choices for the same coordinates.
The matrices $A_N$ and $\Phi_N$ are self-adjoint for the column metric:
\[
A_N^TB_N^{-1}=B_N^{-1}A_N,\qquad
\Phi_N^TB_N^{-1}=B_N^{-1}\Phi_N.
\]
For $A_N$ this follows directly from the adjacent binomial ratios; for
$\Phi_N$ it is binomial duality from Section~\ref{sec:1-4}.
If $\phi_j$ is column $j$ of $\Phi_N$, then
\[
\phi_j^TB_N^{-1}\phi_k=\frac{2^N}{\binom Nj}\delta_{jk}.
\]
Thus $2^{-N/2}\sqrt{\binom Nj}\,\phi_j$ are orthonormal in this
weighted column space, not generally in ordinary Euclidean coordinates.
The coordinate change $u\mapsto B_N^{-1/2}u$ turns this metric into
the Euclidean one and yields $U_N$ of \eqref{eq:ency-U-normalization}.
See \cite[Appendix~A]{KocikFeynman} for the source's vector/covector viewpoint.

\section*{Connection to the supplied article sources}
Section~3 (``Krawtchouk matrices and symmetric tensors'') of the supplied
\texttt{KravFeinsKocik.tex}, and the tensor-product and symmetric-tensor
appendices of \texttt{jkhadamard\_Oct\_08\_00.tex}, supply source material
for this chapter. The full tensor space has dimension $2^N$, whereas its
symmetric sector has dimension $N+1$. Orbit sums, orbit averages, and
orthonormal orbit vectors must be distinguished when comparing matrices.
See also Chapter~\ref{ch:9}.

\chapter{Quantum Probability and Tensor Algebra}
\label{ch:9}
\section{Flip operator and quantum random walk}
\label{sec:9-1}
On $(\mathbb R^2)^{\otimes N}$ let $F_r$ flip the $r$th binary factor
and set $\mathcal A_N=\sum_{r=1}^NF_r$. In the binary-string basis this
is the adjacency operator of the $N$-cube. Choosing one factor uniformly
to flip gives the stochastic operator $\mathcal A_N/N$.
The symmetric sector is indexed by Hamming weight $j$ and yields the
Ehrenfest chain.

Basis normalization matters. For the orbit averages
\[
a_j=\binom Nj^{-1}\sum_{|\epsilon|=j}e_\epsilon,
\]
one obtains $\mathcal A_Na_j=(N-j)a_{j+1}+ja_{j-1}$, hence the book's
matrix $A_N$. The unnormalized orbit sums instead give
$(j+1)s_{j+1}+(N-j+1)s_{j-1}$ with endpoint terms omitted.
The single flip $F$ is also a real quantum gate, but a unitary quantum
walk must be specified separately; the stochastic averaging above is
not itself a unitary evolution.
See \cite[Section~7]{KocikFeynman}; notation follows the book conventions.

\section{Krawtchouk matrices as eigenvectors}
\label{sec:9-2}
The local Hadamard identity $FH=HG$ yields
\[
\mathcal A_NH^{\otimes N}=H^{\otimes N}\mathcal G_N,\qquad
\mathcal G_N=\sum_{r=1}^NG_r.
\]
A binary string of weight $j$ is an eigenvector of $\mathcal G_N$ with
eigenvalue $N-2j$. On the orbit-average symmetric sector, this reduces to the master
equation proved in Section~\ref{sec:2-2} and interpreted through
symmetric powers in Section~\ref{sec:8-1}.
On the full cube the multiplicity is $\binom Nj$; in the symmetric sector
each eigenvalue occurs once. This distinction explains why the
$(N+1)$-dimensional Krawtchouk matrix is smaller than the $2^N$-dimensional
Hadamard tensor matrix. Normalizing $H$ by $1/\sqrt2$ gives an orthogonal
tensor transform.
See \cite[Section~7]{KocikFeynman}; notation follows the book conventions.

\section{Trace formulas and MacMahon's theorem}
\label{sec:9-3}
% TODO: Develop this section; add verified references, examples, and exercises.

\noindent\textbf{TRACE FORMULAS AND MACMAHON'S THEOREM NEEDS TO BE DONE.}


\section{Chebyshev polynomials}
\label{sec:9-4}
% TODO: Develop this section; add verified references, examples, and exercises.

\noindent\textbf{CHEBYSHEV POLYNOMIALS NEEDS TO BE DONE.}


\section{Supersymmetry}
\label{sec:9-5}
% TODO: Develop this section; add verified references, examples, and exercises.

\noindent\textbf{SUPERSYMMETRY NEEDS TO BE DONE.}


\section{Twistons}
\label{sec:9-6}
The twistons of \cite[Section~5, pp.~9--10]{KocikFeynman} are a toy model,
not a claim about physical particles. Represent an untwisted band by
$\epsilon=+1$ and a M\"obius band by $\epsilon=-1$. The model's combination
rule multiplies signs: two negative bands give a positive sign, while one
negative and one positive give a negative sign. This stipulates parity of
twisting; it does not specify a general topological gluing operation.

For $N$ bands, of which $j$ are M\"obius bands, define the order-$n$
interaction energy to be the sum of products over all $n$-element subsets.
By \eqref{eq:sign-energy-Krawtchouk},
\[
E_n=K_n(j;N)=(\Phi_N)_{nj},\qquad 0\leq n\leq N.
\]
Column $j$ of $\Phi_N$ therefore lists all interaction energies of that
configuration. The algebraic identity and its proof belong to
Section~\ref{sec:4-1}; no second polynomial normalization is introduced here.
For $N=3$ and one M\"obius band, the signs $(+1,+1,-1)$ yield
\[
(E_0,E_1,E_2,E_3)=(1,1,-1,-1),
\]
which is column $1$ of $\Phi_3$ in Appendix~\ref{app:matrix-tables}.

\begin{remark}[Counting convention in the source]
Section~5 initially counts nonorientable bands, but its concluding
sentence on page~10 says ``untwisted bands.'' Here $j$ consistently counts
M\"obius bands, as required by \eqref{eq:gen}. Pair interactions mean
$\sum_{a<b}\epsilon_a\epsilon_b$, not an unrestricted double sum.
\end{remark}


\section{Hadamard compression in the Encyclopedia normalization}
\label{sec:uploaded-hadamard-compression}
The supplied \texttt{jkhadamard\_Oct\_08\_00.tex}, Sections~4 and~6,
and \texttt{KravFeinsKocik.tex}, Sections~2--3, connect Hadamard matrices,
Krawtchouk matrices, and symmetric tensors. Here is the compression
formula in the notation fixed by this book.
Set
\begin{equation*}
H=\begin{pmatrix}1&1\\1&-1\end{pmatrix},\qquad
T_N=H^{\otimes N},\qquad
s_j=\sum_{|\epsilon|=j}e_\epsilon,\qquad
b_j=\binom Nj,\qquad a_j=s_j/b_j.
\end{equation*}
For an input string of weight $j$, the sum of the output coefficients
at weight $n$ is the coefficient of $v^n$ in
$(1+v)^{N-j}(1-v)^j$, namely $K_n(j;N)$.
Permutation symmetry then shows that each output coefficient at that
weight in $T_Na_j$ is $K_n(j;N)/b_n$. Consequently
\begin{equation*}
T_Na_j=\sum_{n=0}^N K_n(j;N)a_n.
\end{equation*}
Thus the restriction of $T_N$ to the symmetric sector has matrix
$\Phi_N$ in the orbit-average basis, not in every choice of basis.
In the orthonormal basis $q_j=s_j/\sqrt{b_j}$, the normalized transform
$2^{-N/2}T_N$ has matrix
\begin{equation*}
2^{-N/2}B_N^{-1/2}\Phi_N B_N^{1/2}.
\end{equation*}
This matrix is orthogonal because $H/\sqrt2$ is orthogonal and the
symmetric sector is invariant. This is a unitary (real orthogonal)
transform, distinct from the stochastic flip average $A_N/N$.
The uploaded articles are preserved as historical source documents;
this derivation explicitly fixes the basis and normalization rather
than importing their standalone preambles or macros.

\chapter{Heisenberg Algebra}
\label{ch:10}
\section{Representations of the Heisenberg algebra}
\label{sec:10-1}
% TODO: Develop this section; add verified references, examples, and exercises.

\noindent\textbf{REPRESENTATIONS OF THE HEISENBERG ALGEBRA NEEDS TO BE DONE.}


\section{Raising and velocity operators; number operator}
\label{sec:10-2}
% TODO: Develop this section; add verified references, examples, and exercises.

\noindent\textbf{RAISING AND VELOCITY OPERATORS; NUMBER OPERATOR NEEDS TO BE DONE.}


\section{Evolution structure and Hamiltonian}
\label{sec:10-3}
% TODO: Develop this section; add verified references, examples, and exercises.

\noindent\textbf{EVOLUTION STRUCTURE AND HAMILTONIAN NEEDS TO BE DONE.}


\section{Time-zero polynomials}
\label{sec:10-4}
% TODO: Develop this section; add verified references, examples, and exercises.

\noindent\textbf{TIME-ZERO POLYNOMIALS NEEDS TO BE DONE.}



\chapter{Central Limit Theorem}
\label{ch:11}
\section{Hermite polynomials}
\label{sec:11-1}
The centered generating function has the exact form
\[
G_N(x,v)=(1-v^2)^{N/2}\exp(x\operatorname{artanh}v).
\]
For $x=\sqrt N\xi$ and $v=t/\sqrt N$, its logarithm tends to $\xi t-t^2/2$. Consequently, for fixed $n$,
\[
N^{-n/2}k_n(\sqrt N\xi,N)\longrightarrow\operatorname{He}_n(\xi),
\]
where $\exp(\xi t-t^2/2)=\sum_n\operatorname{He}_n(\xi)t^n/n!$. When polynomial values are interpreted only on the original lattice, take admissible sequences approaching $\xi$.

\section{Discrete stochastic differential equations}
\label{sec:11-2}
The exact finite-step precursor of Hermite martingales is \eqref{eq:discrete-iterated-recursion}. In factorial normalization it reads
\[
k_n(X_{m+1},m+1)-k_n(X_m,m)
=n\,k_{n-1}(X_m,m)\epsilon_{m+1}.
\]
In particular, $X_m$, $X_m^2-m$, and $X_m^3-(3m-2)X_m$ are martingales. The update changes both the spatial variable and the order parameter; it is not the fixed-$N$ urn-chain generator of Chapter~\ref{ch:3}.

To identify the scaling, define polynomials $H_n(x,t)$ by
\[
\exp(ux-tu^2/2)=\sum_{n=0}^{\infty}H_n(x,t)\frac{u^n}{n!}.
\]
For $t>0$, $H_n(x,t)=t^{n/2}\operatorname{He}_n(x/\sqrt t)$, and coefficient differentiation gives
\[
\partial_xH_n=nH_{n-1},\qquad
\partial_tH_n+\tfrac12\partial_x^2H_n=0.
\]
For standard Brownian motion $B_t$, Ito's formula therefore yields
\[
dH_n(B_t,t)=nH_{n-1}(B_t,t)\,dB_t.
\]
This is the continuous counterpart of the displayed finite-step equation. The generating-function expansion proves fixed-degree polynomial scaling; a process-level convergence statement additionally requires a functional central limit argument and appropriate integrability control, neither of which is inferred solely from the formal similarity of the equations.

\section{Bilinear kernels and the Mehler limit}
\label{sec:11-3}
% TODO: Develop this section; add verified references, examples, and exercises.

\noindent\textbf{BILINEAR KERNELS AND THE MEHLER LIMIT NEEDS TO BE DONE.}



\chapter{Clebsch--Gordan Coefficients}
\label{ch:12}
\section{Clebsch--Gordan coefficients and Krawtchouk polynomials}
\label{sec:12-1}
% TODO: Develop this section; add verified references, examples, and exercises.

\noindent\textbf{CLEBSCH--GORDAN COEFFICIENTS AND KRAWTCHOUK POLYNOMIALS NEEDS TO BE DONE.}


\section{Racah coefficients}
\label{sec:12-2}
% TODO: Develop this section; add verified references, examples, and exercises.

\noindent\textbf{RACAH COEFFICIENTS NEEDS TO BE DONE.}


\section{Clebsch--Gordan coefficients for Heisenberg algebras}
\label{sec:12-3}
% TODO: Develop this section; add verified references, examples, and exercises.

\noindent\textbf{CLEBSCH--GORDAN COEFFICIENTS FOR HEISENBERG ALGEBRAS NEEDS TO BE DONE.}



\chapter{Orthogonal Polynomials}
\label{ch:13}
\section{Three-term recurrence in terms of $A$, $K$, and $\Lambda$}
\label{sec:13-1}
Differentiating the generating function gives
\[
k_{n+1}(x,N)=xk_n(x,N)-n(N-n+1)k_{n-1}(x,N).
\]
The spatial equation is
\[
\frac{N-x}{2}\bigl[k_n(x+2,N)-k_n(x,N)\bigr]
+\frac{N+x}{2}\bigl[k_n(x-2,N)-k_n(x,N)\bigr]
=-2n k_n(x,N).
\]
These formulas fix the normalization for later operator treatments.

\section{Nonsymmetric case}
\label{sec:13-2}
% TODO: Develop this section; add verified references, examples, and exercises.


\chapter{Krawtchouk Transforms}
\label{ch:14}
\section{Orthogonal transformation associated to $K$}
\label{sec:14-1}
% TODO: Develop this section; add verified references, examples, and exercises.

\section{Exponential function in the Krawtchouk basis}
\label{sec:14-2}
% TODO: Develop this section; add verified references, examples, and exercises.

\section{Krawtchouk transform}
\label{sec:14-3}
% TODO: Develop this section; add verified references, examples, and exercises.

\section{Krawtchouk transform and FFT}
\label{sec:14-4}
% TODO: Develop this section; add verified references, examples, and exercises.

\section{Binomial bases and dual transforms}
\label{sec:14-5}
Let $P_N=(\binom mn)_{n,m=0}^N$, $J$ be coordinate reversal, and $D_N=\operatorname{diag}(2^m)$. For the binomial column $b_m(n)=\binom mn$,
\[
\Phi_Nb_m=2^m b_{N-m},\qquad P_N^{-1}\Phi_NP_N=JD_N.
\]
This is a nonorthogonal basis. Rayleigh quotients in these coordinates must use the Gram matrix $P_N^TP_N$.


\chapter{Hypergeometric Functions}
\label{ch:15}
\section{Krawtchouk polynomials as hypergeometric functions}
\label{sec:15-1}
% TODO: Develop this section; add verified references, examples, and exercises.

\noindent\textbf{KRAWTCHOUK POLYNOMIALS AS HYPERGEOMETRIC FUNCTIONS NEEDS TO BE DONE.}


\section{Addition formulas}
\label{sec:15-2}
% TODO: Develop this section; add verified references, examples, and exercises.

\noindent\textbf{ADDITION FORMULAS NEEDS TO BE DONE.}



\chapter{Symmetric Krawtchouk Matrices}
\label{ch:16}
\section{The matrix $T$}
\label{sec:16-1}
% TODO: Develop this section; add verified references, examples, and exercises.

\section{$S^2$ and trace formulas}
\label{sec:16-2}
Binomial duality yields the exact bilinear identity
\[
(S_N^2)_{jk}=\binom Nj\binom Nk\sum_{n=0}^NK_n(j;N)K_n(k;N).
\]
Parseval applied to the generating functions evaluates it as
\begin{equation}\label{eq:beta}
(S_N^2)_{jk}=\frac{4^N}{\pi}\binom Nj\binom Nk
\cos\frac{\pi(j-k)}2\,
\mathrm B\left(\frac{j+k+1}{2},\frac{2N-j-k+1}{2}\right).
\end{equation}
The beta integral follows by factoring $1\pm e^{i\theta}$ into half-angle sine and cosine factors. This identity is exact, not an asymptotic statement.

\section{Spectrum of $S$}
\label{sec:16-3}
% TODO: Develop this section; add verified references, examples, and exercises.

\section{Symmetry sectors and block decompositions}
\label{sec:16-4}
Define $(Jf)_j=f_{N-j}$ and $(Ef)_j=(-1)^jf_j$. Then
\[
JS_N=S_NE,\qquad S_NJ=ES_N,
\]
so $S_N$ commutes with $J+E$. For the square, indices of different parity decouple. Conjugation by $\operatorname{diag}((-1)^{\lfloor j/2\rfloor})$ makes each parity block entrywise positive, by \eqref{eq:beta}.

\section{Beta-function kernel and spectral asymptotics}
\label{sec:16-5}
\begin{remark}[Working asymptotic program, not a completed theorem]
Put $\xi=(N-2j)/\sqrt N$ and $\eta=(N-2k)/\sqrt N$. Compact-window Stirling expansions of \eqref{eq:beta}, after sign removal and normalization by $N/8^N$, lead to the candidate kernel
\[
L(\xi,\eta)=\frac{1}{\sqrt2\pi^{3/2}}
\exp\left[-\frac38(\xi^2+\eta^2)+\frac14\xi\eta\right].
\]
Its Gaussian ground state is $\exp(-\xi^2/(2\sqrt2))$, with eigenvalue $\kappa=(\sqrt8-2)/\pi$. The proposed positive-eigenvalue asymptotic is
\[
\lambda_{\max}(S_N)\sim\sqrt{\sqrt8-2}\,
\frac{(\sqrt8)^N}{\sqrt{\pi N}}.
\]
Global operator convergence and control of the positive spectral branch still require proof. Fixed-degree Hermite limits alone do not justify this conclusion.
\end{remark}


\chapter{Gaussian Quadrature}
\label{ch:17}
\noindent\emph{Source and derivation note.}
The orthogonal system and normalization are those of Sections~\ref{sec:1-2}, \ref{sec:5-2}, and \ref{sec:13-1}. The product-martingale source is the supplied \texttt{sources/KravFeinsKocik.tex}, identified in the source register. The Gaussian rule, terminal factorization, proofs, and numerical examples here are developed from those identities; they are not attributed to an unread item in the supplied literature list. Numerical provenance and reproduction details are in \texttt{sources/INTEGRATION\_gaussian\_summation.md} and \texttt{sources/compute\_gaussian\_krawtchouk.py}.

\section{Zeros of Krawtchouk polynomials}
\label{sec:17-1}
The measure in this chapter is finite and discrete:
\[
\mu_N=2^{-N}\sum_{j=0}^N\binom Nj\delta_j,\qquad
\mathcal S_N[f]=2^{-N}\sum_{j=0}^N\binom Nj f(j).
\]
Thus Gaussian quadrature here compresses a weighted sum, not an ordinary integral. Keep two integers distinct: $N$ specifies the $N+1$ support points, whereas $n$ specifies the degree of the orthogonal polynomial and the number of Gaussian nodes. We take $1\le n\le N$ unless stated otherwise.

Extend $K_n(j;N)$ polynomially off the integer lattice via the recurrence of Section~\ref{sec:1-3}. Its leading coefficient in $j$ is $(-2)^n/n!$. The monic polynomials
\[
p_n(j)=\frac{(-1)^n n!}{2^n}K_n(j;N)
=\frac{(-1)^n}{2^n}k_n(N-2j,N)
\]
therefore satisfy
\begin{equation}\label{eq:quadrature-monic-recurrence}
p_{r+1}(j)=(j-N/2)p_r(j)-\frac{r(N-r+1)}4p_{r-1}(j),\qquad p_0=1.
\end{equation}
Their squared norms under $\mu_N$ are
\begin{equation}\label{eq:quadrature-monic-norm}
h_r=\mathcal S_N[p_r^2]=\frac{(r!)^2}{4^r}\binom Nr
=\frac{r!N!}{4^r(N-r)!}.
\end{equation}
The Gaussian nodes $z_1<\cdots<z_n$ are the zeros of $p_n$, equivalently the eigenvalues of the $n\times n$ real symmetric Jacobi matrix
\begin{equation}\label{eq:quadrature-Jacobi}
\mathcal J_{N,n}=\begin{pmatrix}
N/2&b_1& & &\\
b_1&N/2&b_2& &\\
 &b_2&N/2&\ddots&\\
 & &\ddots&\ddots&b_{n-1}\\
 & & &b_{n-1}&N/2
\end{pmatrix},\qquad b_r=\frac12\sqrt{r(N-r+1)}.
\end{equation}
For $n=1$ this means the scalar matrix $(N/2)$. Positivity of $b_r$ gives simple real eigenvalues; the Rayleigh quotient is the weighted average of $j$ against the square of a nonzero polynomial of degree at most $n-1$. Such a polynomial cannot vanish on all $N$ other support points, so $0<z_a<N$. The zeros of successive degrees strictly interlace. Reflection symmetry gives
$z_{n+1-a}=N-z_a$.
These nodes are generally not integers: polynomial extension, rather than a table of lattice values alone, is essential.

\section{Integral zeros}
\label{sec:17-2}
\paragraph{The trivial terminal zeros.}
Remark~\ref{rem:terminal-integral-zeros} gives a universal family of integral zeros at the terminal degree:
\[
k_{N+1}(x,N)=\prod_{j=0}^{N}(x-(N-2j)).
\]
For even $N=2m$ this is $x\prod_{r=1}^{m}(x^2-4r^2)$; for odd $N=2m+1$ it is $\prod_{r=0}^{m}(x^2-(2r+1)^2)$. Thus all centered spectral points are zeros. For example,
\[
k_6(x,5)=(x^2-1)(x^2-9)(x^2-25).
\]
In the uncentered coordinate $j=(N-x)/2$, this is precisely the monic terminal polynomial $p_{N+1}(j)=\prod_{r=0}^{N}(j-r)$ discussed below. This terminal case is automatic and should be distinguished from the nontrivial question of integral zeros in degrees $1\le n\le N$.

Integer zeros are special cases of these polynomial zeros, not a prerequisite for Gaussian summation. Since
$K_n(N-j;N)=(-1)^nK_n(j;N)$, every odd degree has the central zero $j=N/2$; this is integral exactly when $N$ is even. For degree two,
\[
K_2(j;N)=\frac{(N-2j)^2-N}{2},\qquad
z_\pm=\frac{N\pm\sqrt N}{2}.
\]
For $N\ge2$ these two zeros are integers exactly when $N$ is a perfect square: a square and its square root have the same parity. For example, $N=9$ gives nodes $3$ and $6$, each with Gaussian weight $1/2$. Consequently the two-node rule integrates, in the discrete sense, every cubic exactly under $\mu_9$.
For $N=5$ the degree-three centered polynomial is $k_3(x,5)=x(x^2-13)$, giving $j=(5-\sqrt{13})/2,5/2,(5+\sqrt{13})/2$, none integral.
These examples are not a classification of integral zeros at higher degree.
\par\noindent\textbf{A GENERAL CLASSIFICATION OF INTEGRAL ZEROS NEEDS TO BE DONE.}

\section{Gaussian--Krawtchouk summation}
\label{sec:17-3}
For each Gaussian node define the Lagrange polynomial
\[
\ell_a(j)=\frac{p_n(j)}{(j-z_a)p_n'(z_a)},\qquad
\lambda_a=\mathcal S_N[\ell_a].
\]
Then the $n$-node formula is
\begin{equation}\label{eq:gaussian-Krawtchouk-sum}
2^{-N}\sum_{j=0}^N\binom Nj f(j)
=\sum_{a=1}^n\lambda_a f(z_a),\qquad \deg f\le2n-1.
\end{equation}
To prove exactness, divide $f=q p_n+r$ with $\deg q,\deg r\le n-1$. Orthogonality annihilates $\mathcal S_N[q p_n]$, while interpolation of $r$ evaluates its sum exactly. Since $p_n(z_a)=0$, the right side is also unchanged by removing $q p_n$.
Applying exactness to $\ell_a^2$ shows
$\lambda_a=\mathcal S_N[\ell_a^2]>0$, and applying it to $1$ gives $\sum_a\lambda_a=1$. The rule is symmetric: $\lambda_{n+1-a}=\lambda_a$.
Two equivalent computable weight formulas are
\begin{equation}\label{eq:quadrature-Christoffel}
\lambda_a=\left[\sum_{r=0}^{n-1}\frac{K_r(z_a;N)^2}{\binom Nr}\right]^{-1}
=\frac{h_{n-1}}{p_n'(z_a)p_{n-1}(z_a)}.
\end{equation}
Alternatively, if $u_a$ is the unit eigenvector of \eqref{eq:quadrature-Jacobi} for $z_a$, then $\lambda_a=|(u_a)_0|^2$. The spectral measure of the first coordinate gives this identity; the Christoffel formula follows by summing squares of orthonormal polynomials. These formulas use the probability-normalized binomial measure. For the unnormalized sum $\sum_j\binom Nj f(j)$ multiply every weight by $2^N$.

\subsection{Error and the finite-support endpoint}
For $f\in C^{2n}([0,N])$, Hermite interpolation at the $n$ nodes, followed by summation against the positive measure, gives
\[
\mathcal S_N[f]-\sum_a\lambda_a f(z_a)
=\frac{h_n}{(2n)!}f^{(2n)}(\xi),\qquad \xi\in[0,N].
\]
In particular, the error on the monic polynomial $j^{2n}$ is exactly $h_n>0$, so degree $2n-1$ is sharp. The rule remains meaningful for nonpolynomial functions provided values between lattice points are defined. If only the $N+1$ discrete data values are known, choosing an off-lattice interpolation is extra information; one cannot assume an arbitrary data array has prescribed Gaussian-node values.
For $n=N+1$, the terminal polynomial is $p_{N+1}(j)=\prod_{r=0}^N(j-r)$ and its norm is zero. The full Jacobi matrix recovers the original nodes $0,\ldots,N$ and weights $2^{-N}\binom Nj$. This is the original finite sum, exact for every lattice function; it is not a new compressed rule, and the positive-norm error discussion above is restricted to $n\le N$.

\subsection{Accurate examples for $N=5,10,15$}
The following tables give both reduced-node rules $(N,n)=(5,3),(10,5),(15,8)$ and the degree-$N$ rules $(5,5),(10,10),(15,15)$. Thus the reduced rules replace respectively $6,11,16$ summands by $3,5,8$ evaluations and are exact through degrees $5,9,15$. The degree-$N$ rules are exact through degrees $9,19,29$. All nodes below are in the original $j$ coordinate; centered nodes are $x_a=N-2z_a$ with the same weights.

For $(N,n)=(5,3)$ one also has the exact formula
\[
\mathcal S_5[f]=\frac5{26}f\!\left(\frac{5-\sqrt{13}}2\right)
+\frac8{13}f\!\left(\frac52\right)
+\frac5{26}f\!\left(\frac{5+\sqrt{13}}2\right),\qquad\deg f\le5.
\]
The tables show 18 significant digits. Their machine-readable companions retain 65 digits. Computation used a symmetric Jacobi eigensystem at 90-decimal working precision, independently checked against the centered polynomial recurrence, the Christoffel weight formula, positivity, and every scaled moment $(j/N)^r$ for $0\le r\le2n-1$. Maximum absolute scaled-moment error over all six rules was below $10^{-85}$ at working precision. These numerical residuals are consistency checks, not interval-certified error bounds; the printed rounded tables have correspondingly lower accuracy.
\par\medskip\noindent\textbf{$N=5$, $n=3$; exact through degree $5$.}
\begin{center}\small
\begin{tabular}{rll}
$a$ & $z_a$ & $\lambda_a$\\\hline
1 & 0.697224362268005327 & 0.192307692307692318\\
2 & 2.5 & 0.615384615384615419\\
3 & 4.30277563773199478 & 0.192307692307692318\\
\end{tabular}
\end{center}
\par\medskip\noindent\textbf{$N=5$, $n=5$; exact through degree $9$.}
\begin{center}\small
\begin{tabular}{rll}
$a$ & $z_a$ & $\lambda_a$\\\hline
1 & 0.0649539076702599777 & 0.0406121317521831102\\
2 & 1.24678392596102428 & 0.244622767576675948\\
3 & 2.5 & 0.42953020134228187\\
4 & 3.7532160740389755 & 0.244622767576675948\\
5 & 4.93504609232974012 & 0.0406121317521831102\\
\end{tabular}
\end{center}
\par\medskip\noindent\textbf{$N=10$, $n=5$; exact through degree $9$.}
\begin{center}\small
\begin{tabular}{rll}
$a$ & $z_a$ & $\lambda_a$\\\hline
1 & 1.0 & 0.01953125\\
2 & 3.0 & 0.234375\\
3 & 5.0 & 0.4921875\\
4 & 7.0 & 0.234375\\
5 & 9.0 & 0.01953125\\
\end{tabular}
\end{center}
\par\medskip\noindent\textbf{$N=10$, $n=10$; exact through degree $19$.}
\begin{center}\small
\begin{tabular}{rll}
$a$ & $z_a$ & $\lambda_a$\\\hline
1 & 0.00426140162761967471 & 0.00100005627585995687\\
2 & 1.03215639227186573 & 0.0108643711646248114\\
3 & 2.1094239579268903 & 0.0553578832768971893\\
4 & 3.24085911987510711 & 0.160428693955712715\\
5 & 4.40940159989234548 & 0.272348995326905319\\
6 & 5.59059840010765452 & 0.272348995326905319\\
7 & 6.75914088012489334 & 0.160428693955712715\\
8 & 7.8905760420731097 & 0.0553578832768971893\\
9 & 8.96784360772813471 & 0.0108643711646248114\\
10 & 9.99573859837238032 & 0.00100005627585995687\\
\end{tabular}
\end{center}
\par\medskip\noindent\textbf{$N=15$, $n=8$; exact through degree $15$.}
\begin{center}\small
\begin{tabular}{rll}
$a$ & $z_a$ & $\lambda_a$\\\hline
1 & 0.810168654648509734 & 0.000551841058012324075\\
2 & 2.70329267453015554 & 0.0178684857525261559\\
3 & 4.61682155336053412 & 0.134507504747883117\\
4 & 6.53821132811295858 & 0.347072168441578399\\
5 & 8.46178867188704054 & 0.347072168441578399\\
6 & 10.3831784466394659 & 0.134507504747883117\\
7 & 12.2967073254698445 & 0.0178684857525261559\\
8 & 14.1898313453514895 & 0.000551841058012324075\\
\end{tabular}
\end{center}
\par\medskip\noindent\textbf{$N=15$, $n=15$; exact through degree $29$.}
\begin{center}\small
\begin{tabular}{rll}
$a$ & $z_a$ & $\lambda_a$\\\hline
1 & 0.000210675294063239482 & 0.0000305592264942427708\\
2 & 1.00264068454641486 & 0.000462988590523494585\\
3 & 2.01486605080593995 & 0.00335579980348639896\\
4 & 3.05052871756057487 & 0.0155659622211499205\\
5 & 4.12061368475752943 & 0.0506438220470849867\\
6 & 5.22535900511131501 & 0.117850851005629698\\
7 & 6.35599416357050373 & 0.195969221661716791\\
8 & 7.5 & 0.232241590887828919\\
9 & 8.64400583642949627 & 0.195969221661716791\\
10 & 9.77464099488868499 & 0.117850851005629698\\
11 & 10.8793863152424706 & 0.0506438220470849867\\
12 & 11.9494712824394256 & 0.0155659622211499205\\
13 & 12.9851339491940596 & 0.00335579980348639896\\
14 & 13.9973593154535845 & 0.000462988590523494585\\
15 & 14.9997893247059366 & 0.0000305592264942427708\\
\end{tabular}
\end{center}


\chapter{Coding Theory}
\label{ch:18}
\section{MacWilliams's theorem}
\label{sec:18-1}
% TODO: Develop this section; add verified references, examples, and exercises.

\noindent\textbf{MACWILLIAMS'S THEOREM NEEDS TO BE DONE.}


\section{Association schemes}
\label{sec:18-2}
% TODO: Develop this section; add verified references, examples, and exercises.

\noindent\textbf{ASSOCIATION SCHEMES NEEDS TO BE DONE.}



\appendix
\chapter{Tables of Krawtchouk and Symmetric Matrices}
\section{K and S matrices for orders 1 through 14}
% TODO: Populate and verify.


\chapter{Polynomial Tables}
\label{app:polynomial-tables}
The developed tables cover $N=3,5$, with every degree $0\leq n\leq N$.
The normalization is
\[
(1+v)^{N-j}(1-v)^j=\sum_{n=0}^N K_n(j;N)v^n,
\qquad k_n(N-2j,N)=n!K_n(j;N).
\]
All lattice entries have been verified exactly against the binomial
coefficient formula and Appendix~\ref{app:matrix-tables}.

\section{Centered variables $x$ and $N$}
From $k_0=1$, $k_1=x$ and
$k_{n+1}=xk_n-n(N-n+1)k_{n-1}$, we obtain
\begin{align*}
k_{0}(x,N)&=1 \\
k_{1}(x,N)&=x \\
k_{2}(x,N)&=- N + x^{2} \\
k_{3}(x,N)&=- 3 N x + x^{3} + 2 x \\
k_{4}(x,N)&=3 N^{2} - 6 N x^{2} - 6 N + x^{4} + 8 x^{2} \\
k_{5}(x,N)&=15 N^{2} x - 10 N x^{3} - 50 N x + x^{5} + 20 x^{3} + 24 x
\end{align*}
\subsection{$N=3$}
The lattice is $x=3,1,-1,-3$.
\begin{align*}
k_{0}(x,3)&=1 \\
k_{1}(x,3)&=x \\
k_{2}(x,3)&=x^{2} - 3 \\
k_{3}(x,3)&=x^{3} - 7 x
\end{align*}
\subsection{$N=5$}
The lattice is $x=5,3,1,-1,-3,-5$.
\begin{align*}
k_{0}(x,5)&=1 \\
k_{1}(x,5)&=x \\
k_{2}(x,5)&=x^{2} - 5 \\
k_{3}(x,5)&=x^{3} - 13 x \\
k_{4}(x,5)&=x^{4} - 22 x^{2} + 45 \\
k_{5}(x,5)&=x^{5} - 30 x^{3} + 149 x
\end{align*}
\section{Degree and lattice variables $i$ and $j$}
Rows are degrees $i=0,\ldots,N$ and columns are lattice indices
$j=0,\ldots,N$. These are $\Phi_N$ matrices of $K_i(j;N)$ values,
not factorial-scaled $k_i$ values.
\begingroup
\small
\setlength{\arraycolsep}{3pt}
\[
\bigl(K_i(j;3)\bigr)_{i,j=0}^{3}=\begin{pmatrix}
1 & 1 & 1 & 1 \\
3 & 1 & -1 & -3 \\
3 & -1 & -1 & 3 \\
1 & -1 & 1 & -1
\end{pmatrix}.
\]
\endgroup
\begingroup
\small
\setlength{\arraycolsep}{3pt}
\[
\bigl(K_i(j;5)\bigr)_{i,j=0}^{5}=\begin{pmatrix}
1 & 1 & 1 & 1 & 1 & 1 \\
5 & 3 & 1 & -1 & -3 & -5 \\
10 & 2 & -2 & -2 & 2 & 10 \\
10 & -2 & -2 & 2 & 2 & -10 \\
5 & -3 & 1 & 1 & -3 & 5 \\
1 & -1 & 1 & -1 & 1 & -1
\end{pmatrix}.
\]
\endgroup
\section{Lattice variables $j$ and $N$; orders 1 through 20}
The following formulas supply degrees zero through five. For fixed $N$,
retain only $n\leq N$. The orders completed here are $N=3,5$;
the other orders in the original program remain to be developed.
\begin{align*}
K_{0}(j;N)&=1 \\
K_{1}(j;N)&=N - 2 j \\
K_{2}(j;N)&=\frac{N^{2} - 4 N j - N + 4 j^{2}}{2} \\
K_{3}(j;N)&=- \frac{\left(- N + 2 j\right) \left(N^{2} - 4 N j - 3 N + 4 j^{2} + 2\right)}{6} \\
K_{4}(j;N)&=\frac{N^{4} - 8 N^{3} j - 6 N^{3} + 24 N^{2} j^{2} + 24 N^{2} j + 11 N^{2} - 32 N j^{3} - 24 N j^{2} - 32 N j - 6 N + 16 j^{4} + 32 j^{2}}{24} \\
K_{5}(j;N)&=- \frac{\left(- N + 2 j\right) \left(N^{4} - 8 N^{3} j - 10 N^{3} + 24 N^{2} j^{2} + 40 N^{2} j + 35 N^{2} - 32 N j^{3} - 40 N j^{2} - 80 N j - 50 N + 16 j^{4} + 80 j^{2} + 24\right)}{120}
\end{align*}
\subsection{$N=3$}
\begin{align*}
K_{0}(j;3)&=1 \\
K_{1}(j;3)&=3 - 2 j \\
K_{2}(j;3)&=2 j^{2} - 6 j + 3 \\
K_{3}(j;3)&=- \frac{\left(2 j - 3\right) \left(2 j^{2} - 6 j + 1\right)}{3}
\end{align*}
\subsection{$N=5$}
\begin{align*}
K_{0}(j;5)&=1 \\
K_{1}(j;5)&=5 - 2 j \\
K_{2}(j;5)&=2 \left(j^{2} - 5 j + 5\right) \\
K_{3}(j;5)&=- \frac{2 \left(2 j - 5\right) \left(j^{2} - 5 j + 3\right)}{3} \\
K_{4}(j;5)&=\frac{2 j^{4} - 20 j^{3} + 64 j^{2} - 70 j + 15}{3} \\
K_{5}(j;5)&=- \frac{\left(2 j - 5\right) \left(2 j^{4} - 20 j^{3} + 60 j^{2} - 50 j + 3\right)}{15}
\end{align*}

\chapter{Spectral Tables}
\section{Eigenvalues of $S$}
% TODO: Populate and verify.

\noindent\textbf{EIGENVALUES OF $S$ NEEDS TO BE DONE.}


\section{Normalization and numerical verification}
% TODO: Populate and verify.

\noindent\textbf{NORMALIZATION AND NUMERICAL VERIFICATION NEEDS TO BE DONE.}



\chapter{Multivariate Extensions}
\section{Remarks on the multivariate case}
% TODO: Populate and verify.

\noindent\textbf{REMARKS ON THE MULTIVARIATE CASE NEEDS TO BE DONE.}



\chapter{Time-Zero Polynomial Tables}
\section{Definitions and tables}
% TODO: Populate and verify.

\noindent\textbf{DEFINITIONS AND TABLES NEEDS TO BE DONE.}



\chapter{Mikhail Kravchuk: A Biographical Sketch}
\section{Biographical chronology}
% TODO: Populate and verify.

\noindent\textbf{BIOGRAPHICAL CHRONOLOGY NEEDS TO BE DONE.}


\section{Primary sources and historical bibliography}
% TODO: Populate and verify.

\noindent\textbf{PRIMARY SOURCES AND HISTORICAL BIBLIOGRAPHY NEEDS TO BE DONE.}



\backmatter
\chapter*{Bibliography and Source Register}
\addcontentsline{toc}{chapter}{Bibliography and Source Register}
\section*{Supplied source materials}
\begin{itemize}
\item \texttt{KravFeinsKocik.tex}: the article source supplying the encyclopedia outline in Section 6.1.
\item \texttt{Kravchuk.zip}: the archived encyclopedia website supplying additional headings and the introductory chapter.
\item \texttt{KravConvolutions.pdf}: supplied supporting material for future transform and convolution chapters.
\end{itemize}
\section*{Literature bibliography}
% TODO: Insert a verified bibliography, then cite it in the developed chapters.
No external bibliographic records have been asserted or reconstructed in this initial skeleton. Before publication, verify authors, titles, dates, venues, and identifiers for every reference.
\end{document}
